\documentclass[UTF8,a4paper,11pt]{ctexart}
\usepackage[margin=2.5cm]{geometry}
\usepackage{booktabs,tabularx,url}
\usepackage[hidelinks]{hyperref}
\setlength{\parindent}{2em}
\title{AI工具使用详情}
\author{}
\date{2026年9月10日}
\begin{document}
\maketitle

\section{工具与使用范围}
本任务使用OpenAI Codex桌面应用及其辅助代理、代码执行、文件处理和文献检索功能。模型名称以实际会话或客户端显示为准，本文档不编造未核验的内部子版本。先前ChatGPT对话“母模型建模思路”和“论文收集与匹配”作为建模路线及检索线索使用；其中数学表达和书目内容并未直接视为事实，而是与赛题、原始附件及可获得的一手文献信息核对。

AI参与了统一建模方案整理、预测与优化程序编写、数据计算、场景与消融实验设计、图表生成、书目信息核验、论文初稿撰写和附件5数据填报。其参与范围包含实质性的模型及代码内容，并不限于语言润色。

本轮修订还使用辅助代理分别处理计算对照、时间映射与工作簿导出等独立工作，主流程汇总结果。代理输出以实际保存的代码、数组和核对记录为依据，不把模型生成的结论直接当作已验证事实。

\section{主要提示与过程}
用户要求结合上述两段对话，遵循C题，在指定LaTeX模板目录完成论文并填写附件5，保留原始代码，同时明确由用户后续编译与审查。过程中用户进一步要求保留附件5原始时间标签，从0:10--0:20开始，以次日0:00--0:10结束。上述约束被纳入数据导出及论文的时间映射说明。

工作过程依次包括：读取赛题与模板，提取附件1至4；按可获得信息选择历史预测，保留日内相关残差；建立线性规划与当前时段反馈执行器；从1月初态连续推进全年；生成指定日期、全年策略和对照实验；依据真实计算结果撰写正文，输出原模板工作簿与源代码。原始附件及模板代码备份另存，未用伪造实测结果填表。

随后用户提供详细审阅意见，要求将重点调整为结算机制驱动的采购决策及规划与实际控制的一致性，而非预先认定自适应风险项更优。用户同时反馈其独立复算问题一得到购电费用35126.9486元、购电量59482.6990\,kWh。本文件只记录这一已提供的复算反馈，不将其扩大为用户已独立复现问题二至四或已核验修订后的全部结果。

\section{本轮实际修订范围}
修订程序以真实结算费用选择经济主案，保留五倍紧急购电与1.5倍上调规则。新增单时段80\%采购分位数、给定原计划下70\%上调阈值及随机价格加权分位数的推导，并注明其不直接适用于完整跨时段储能策略。免费弃余下充放电互斥的线性规划处理保留。

预测和控制配置使用1月15日至31日的实际结算费用验证，保存于\texttt{data/selection\_v2.json}，在2至12月评价前冻结。各问题沿用其原版已执行的1月前缀，同一问题的主案及对照具有共同初态，不用事后选出的配置回写1月。当前冻结主案均采用28个历史场景、无附加CVaR与备用、规则储能执行；每十分钟的内层储能LP和合法跨日预报续接已编程实现并作为对照，未将其描述为附件主案或无条件更优。

新增诊断固定同一购电合同、同一储电初态与同一历史场景，只比较真实结算项目，分别记录规划与因果回放的费用和SOC差异。内层LP回放逐时揭示当前场景值，未来部分保持在当时合法预测，不把完整未来场景提前交给控制器。滚动诊断窗口重叠，其均值和差异不能直接作为全年节费或全局最优差距。

本轮完成23个命名对照和2个主策略引用，共25个完成条目；相同配置复用结果，条目数不代表独立仿真次数。对照覆盖28与12场景、场景期望与点预测、备用单独作用、固定与自适应风险/备用、一月匹配的平均保守程度、仅零时计划、状态/光伏/负荷反馈分解、规则控制、内层LP、跨日价值和退款解释。配置及数值保存在\texttt{data/experiments\_v2.json}及配套数组中，2至12月比较结果不回写一月主案选择。原版偏差与内层LP回放两项诊断均已完成；后者包含四个指定日期的16个决策点，两种控制器各回放448条场景轨迹。显式计入未来修订机会的非预见性场景树未在本轮实现。

时间处理增加统一的144行映射，明确右端点平均功率是建模解释，而不是题目明示事实。原购电模板从0:10--0:20到次日0:00--0:10的标签逐字保留；论文物理区间、四小时汇总和紧急区间从同一索引生成，计算数组没有按模板标签平移。原始代码、模板与修订前结果分别保留，附件由同一新版结果JSON生成。

\section{按问题重排与图形补充}

用户进一步要求一个问题对应一个大标题，并在各问下依次编排问题分析、模型建立与求解。AI据此重组四问正文，展开变量、目标函数和约束含义，将相关机制推导与结果解释放回对应问题。重排前源稿与代码另存\texttt{versions/v2\_before\_question\_restructure/}，此前确定的执行数组与五个工作簿继续保留。

新增热力矩阵、能量桑基图、费用瀑布图及场景SOC阴影带均从已保存数据计算。另对334个历史窗口实际运行12代表日k-medoids缩减，记录每轮几何距离目标，绘制中位数与P10--P90收敛带；这不是主LP求解器的迭代轨迹，也不声称是经济费用的收敛证明。SOC阴影带表示历史场景经验分布，未将其标为统计置信区间。图形保留独立Python源程序、原始曲线及来源说明；程序生成的矢量PDF仅为插图，不是LaTeX编译结果。

\section{已实施的计算核对}
程序核对了数据维度、缺失与非负性、部分预测的未来信息扰动不变性、逐时供需平衡、储电量边界、功率上限、充放电互斥、跨日状态连续、费用分解，以及结果工作簿与求解数组的一致性。文献书目使用DOI注册信息及出版社入口核验，并明确了摘要或可用全文不能支持的细节。以上核对由计算程序及AI协助完成，不等于参赛人员已经逐项人工确认。

修订版进一步检查了两时段储能未来价值行为例，以及改变未观测实值和未发布预报后已执行动作前缀是否不变。新版五份工作簿的16个工作表已预览，计划与调整区域共288720个数值逐格核对；原标签、334天日期、四小时聚合、紧急时段合并和Excel合计缓存一致。文件清单记录实际路径、SHA-256及结果JSON哈希，用于区分已填写结果与空白原模板。这些是有限的数值及文件检查，不等于所有输入上的形式证明，也不替代人工论文审查。

\section{采纳、人工修改与核验状态}
本次交付内容为AI协助形成的论文源码、计算源程序与结果工作簿。按用户要求，本次没有安装LaTeX、没有编译论文，也没有生成论文审稿报告。用户尚未完成或未在本会话中记录的人工修改和人工核验，不在此声称已完成。

本轮根据用户审阅意见修改正文与程序，属于修订工作。目录中保留的原\texttt{example.pdf}或旧\texttt{main.pdf}不是本轮源码的编译证明，程序生成的PDF插图也不是论文成稿。新增复现入口仅作PowerShell语法检查；各底层步骤分别执行，没有为入口测试重复整套年度计算。

用户在最终采用前，应根据实际完成情况补充：参与人工核验人员及日期、模型和结算假设是否采纳、时间模板映射是否确认、主要代码与数值的人工复核记录，以及编译排版后的修改。补充内容应反映真实操作，不能将计划进行的核验写成已经完成。

\section{对应材料}
论文入口为\texttt{main.tex}；新版复现入口为\texttt{code/run\_revision.ps1}，原\texttt{code/run\_all.ps1}保留；全部计算源文件位于\texttt{code/}；原模板位于\texttt{originals/}；修订前文件位于\texttt{versions/v1\_before\_settlement\_revision/}；原始数据摘要、冻结选择、实验及逐时结果位于\texttt{data/}；五个结果工作簿保存于题目指定的附件5目录。主数据接口为\texttt{data/results\_export.json}，版本为\texttt{v2\_settlement\_control}。本文件由用户后续单独编译为AI工具使用详情PDF。
\end{document}
